% =====================================================================
%  2. Contexto del esfuerzo de prueba   (actividad TP1 de ISO/IEC/IEEE 29119-2)
% =====================================================================
\section{Contexto}

\begin{notanormativa}[Correspondencia con el modelo de procesos]
Este apartado desarrolla la actividad \textbf{TP1 «Comprender el contexto»} de \normap{2}:
identificar partes interesadas, analizar el elemento de prueba y la base de prueba, consultar
la política y estrategia aplicables y registrar las restricciones de tiempo, recursos y
regulación.
\end{notanormativa}

% ---------------------------------------------------------------------
\subsection{Proyecto caso de estudio}

\begin{instruccion}
Dé antecedentes del proyecto a probar: autores originales, objetivos, alcances y límites,
plataformas en que está desarrollado, funcionalidades generales que aborda y cualquier otra
información pertinente. Debe ser lo bastante explícito para que un lector externo entienda
qué hace el sistema sin consultar otros documentos.
\end{instruccion}

\campo{Describa aquí el proyecto caso de estudio}

% ---------------------------------------------------------------------
\subsection{Objetivo de la prueba}

\begin{instruccion}
Indique el objetivo general \textbf{de la prueba}, no del sistema. Un objetivo de prueba
enuncia qué se pretende averiguar o demostrar y para qué decisión sirve; no describe las
funciones del software.

Formule también objetivos específicos, cada uno asociado a los riesgos que pretende reducir
(apartado~\ref{sec:riesgos}). Esa asociación es la que permite justificar después la estrategia.
\end{instruccion}

\begin{ejemplo}
\textbf{Objetivo de prueba (correcto):} «Determinar si el módulo de facturación calcula
correctamente totales, impuestos y descuentos bajo las combinaciones de reglas de negocio
vigentes, con el fin de sustentar la decisión de liberación.»

\textbf{Objetivo del sistema (incorrecto en este apartado):} «Permitir a los usuarios emitir
facturas electrónicas.»
\end{ejemplo}

\campo{Objetivo general de la prueba}

\begin{table}[H]
\centering
\caption{Objetivos específicos de la prueba y riesgos que pretenden reducir.}
\label{tab:objetivos}
\begin{tabularx}{\textwidth}{|C{1.6cm}|Y|C{2.2cm}|C{3.2cm}|}
  \hline
  \filaenc \textbf{Clave} & \textbf{Objetivo específico de la prueba} & \textbf{Riesgo(s) asociado(s)} & \textbf{Cómo se evidenciará} \\ \hline
  \campo{OBJ-01} & \campo{Qué se pretende averiguar} & \campo{RP-01} & \campo{Registro o métrica} \\ \hline
  & & & \\ \hline
  & & & \\ \hline
\end{tabularx}
\end{table}

El Cuadro~\ref{tab:objetivos} vincula cada objetivo específico con los riesgos del
apartado~\ref{sec:riesgos} y con la evidencia que permitirá comprobar si el objetivo se
alcanzó; sin esa columna, el objetivo no es verificable al cierre.

% ---------------------------------------------------------------------
\subsection{Objeto de prueba y base de prueba}
\label{sec:objeto-base}

\begin{notanormativa}[Corrección respecto de la plantilla original]
La plantilla original equiparaba los «elementos de prueba» con los \emph{casos de uso}. Son
cosas distintas y conviene registrarlas por separado:
\begin{itemize}
  \item El \textbf{objeto de prueba} es lo que se somete a prueba: sistema, subsistema,
        módulo, componente, servicio, API o interfaz.
  \item La \textbf{base de prueba} es la fuente contra la que se juzga ese objeto: requisitos,
        historias de usuario, casos de uso, reglas de negocio, arquitectura o diseño.
\end{itemize}
Un caso de uso es, casi siempre, base de prueba y no objeto de prueba. Separarlos permite
aplicar esta plantilla a proyectos que no estén documentados mediante casos de uso.
\end{notanormativa}

\begin{instruccion}
Identifique primero los objetos bajo prueba con su versión exacta, y después la base de prueba
que se usará para derivar las condiciones de prueba. Si alguna parte del sistema carece de base
documentada, dígalo: es una restricción real que condiciona las técnicas aplicables
(apartado~\ref{sec:tecnicas}) y suele constituir un riesgo de proyecto.
\end{instruccion}

\begin{table}[H]
\centering
\caption{Objetos bajo prueba, con su versión y criticidad asignada.}
\label{tab:objetos}
\begin{tabularx}{\textwidth}{|C{1.6cm}|L{3.4cm}|C{2.1cm}|C{1.9cm}|Y|}
  \hline
  \filaenc \textbf{Clave} & \textbf{Objeto de prueba} & \textbf{Tipo} & \textbf{Versión} & \textbf{Criticidad asignada (\ref{sec:criticidad})} \\ \hline
  \campo{OP-01} & \campo{Módulo de facturación} & \campo{Módulo} & \campo{v1.2.0} & \campo{Crítico} \\ \hline
  \campo{OP-02} & \campo{API de autenticación} & \campo{Servicio} & \campo{v1.2.0} & \campo{Catastrófico} \\ \hline
  & & & & \\ \hline
\end{tabularx}
\notatabla{Tipo sugerido: sistema · subsistema · módulo · componente · servicio · API · interfaz de usuario · interfaz externa · base de datos.}
\end{table}

\begin{table}[H]
\centering
\caption{Base de prueba utilizada para derivar las condiciones de prueba.}
\label{tab:base-prueba}
\begin{tabularx}{\textwidth}{|C{1.6cm}|L{3.9cm}|C{2.3cm}|C{2.0cm}|Y|}
  \hline
  \filaenc \textbf{Clave} & \textbf{Elemento de la base} & \textbf{Tipo} & \textbf{Referente (\ref{tab:referentes})} & \textbf{Objeto(s) que sustenta} \\ \hline
  \campo{BP-01} & \campo{RF-12 Cálculo de total} & \campo{Requisito} & \campo{REF-02} & \campo{OP-01} \\ \hline
  \campo{BP-02} & \campo{CU-03 Emitir factura} & \campo{Caso de uso} & \campo{REF-02} & \campo{OP-01} \\ \hline
  & & & & \\ \hline
\end{tabularx}
\notatabla{Tipo sugerido: requisito funcional · requisito no funcional · historia de usuario · caso de uso · regla de negocio · documento de arquitectura · especificación de interfaz · norma aplicable.}
\end{table}

Los Cuadros~\ref{tab:objetos} y~\ref{tab:base-prueba} responden a dos preguntas distintas:
\emph{qué se prueba} y \emph{contra qué se juzga}. Sus claves son el punto de partida de la
matriz de trazabilidad del \ref{anx:trazabilidad}, que encadena base de prueba, riesgo,
condición, caso, ejecución, evidencia e incidente.

% ---------------------------------------------------------------------
\subsection{Criticidad asignada a los objetos de prueba}
\label{sec:criticidad}

\begin{politicaacademica}[Política académica: escala de criticidad (SIL)]
El curso emplea una escala de cuatro niveles para clasificar la criticidad de los objetos de
prueba: \textbf{Catastrófico}, \textbf{Crítico}, \textbf{Marginal} y \textbf{Despreciable}.
Esta escala es una \emph{regla local de la Experiencia Educativa}. Describa en el
Cuadro~\ref{tab:sil} el significado que su equipo da a cada nivel, quién lo asigna y con qué
criterio.
\end{politicaacademica}

\begin{notanormativa}[Por qué esta escala no se atribuye a \norma]
La revisión señaló que las etiquetas anteriores se presentaban como si procedieran de \norma.
No es así: \norma\ no define esa escala de cuatro niveles. El marco internacional para niveles
de integridad es \textbf{ISO/IEC/IEEE 15026-3} (\emph{Systems and software assurance —
Integrity levels}), que exige definir los niveles, los requisitos asociados a cada uno y el
método de asignación. Mientras esos elementos no se definan y verifiquen contra dicha norma,
la escala debe tratarse como una clasificación académica interna y no como un nivel de
integridad en sentido normativo.
\end{notanormativa}

\begin{instruccion}
Complete el cuadro siguiente antes de usar la escala. Un nivel de criticidad sin definición ni
método de asignación no puede justificar decisiones de cobertura ni de técnicas.
\end{instruccion}

\begin{table}[H]
\centering
\caption{Definición local de la escala de criticidad y su efecto sobre la estrategia.}
\label{tab:sil}
\begin{tabularx}{\textwidth}{|C{2.3cm}|Y|C{2.6cm}|L{3.5cm}|}
  \hline
  \filaenc \textbf{Nivel} & \textbf{Significado acordado por el equipo (consecuencia de un fallo)} & \textbf{Quién lo asigna y cuándo} & \textbf{Efecto sobre la estrategia} \\ \hline
  Catastrófico & \campo{p.\,ej. pérdida económica irreversible, daño legal o de seguridad} & \campo{Rol y momento} & \campo{Cobertura objetivo, técnicas y niveles} \\ \hline
  Crítico      & \campo{Significado} & \campo{Rol y momento} & \campo{Efecto} \\ \hline
  Marginal     & \campo{Significado} & \campo{Rol y momento} & \campo{Efecto} \\ \hline
  Despreciable & \campo{Significado} & \campo{Rol y momento} & \campo{Efecto} \\ \hline
\end{tabularx}
\end{table}

El Cuadro~\ref{tab:sil} hace explícito qué significa cada nivel y cómo modifica la estrategia;
la asignación concreta a cada objeto se registra en el Cuadro~\ref{tab:objetos}. El criterio de
asignación debe ser la \emph{consecuencia de un fallo}, no la dificultad de probar el elemento.

% ---------------------------------------------------------------------
\subsection{Alcance de la prueba}
\label{sec:alcance}

\begin{notanormativa}[Corrección respecto de la plantilla original]
La plantilla original orientaba las características a probar únicamente hacia los
\emph{requisitos no funcionales}. El alcance debe declarar explícitamente ambos tipos:
características \textbf{funcionales} (calcular un total, autorizar una operación) y
\textbf{no funcionales} (rendimiento, seguridad, usabilidad, compatibilidad). Además, toda
exclusión debe justificarse y relacionarse con el riesgo que se asume al excluirla.
\end{notanormativa}

\begin{instruccion}
Llene los tres cuadros siguientes. Para las características no funcionales puede apoyarse en el
modelo de calidad de producto de \textbf{ISO/IEC 25010:2023} (adecuación funcional, eficiencia
de desempeño, compatibilidad, capacidad de interacción, fiabilidad, seguridad, mantenibilidad,
flexibilidad y seguridad física) y en los requisitos no funcionales del propio sistema. Si cita
un anexo de \normap{4}, indique la edición consultada.
\end{instruccion}

\begin{table}[H]
\centering
\caption{Características funcionales incluidas en el alcance.}
\label{tab:alcance-f}
\begin{tabularx}{\textwidth}{|C{1.6cm}|C{1.7cm}|Y|C{2.2cm}|C{2.1cm}|}
  \hline
  \filaenc \textbf{Clave} & \textbf{Objeto (\ref{tab:objetos})} & \textbf{Característica funcional a probar} & \textbf{Base (\ref{tab:base-prueba})} & \textbf{Prioridad} \\ \hline
  \campo{CF-01} & \campo{OP-01} & \campo{Cálculo de total con descuentos} & \campo{BP-01} & \campo{Alta} \\ \hline
  & & & & \\ \hline
  & & & & \\ \hline
\end{tabularx}
\end{table}

\begin{table}[H]
\centering
\caption{Características no funcionales incluidas en el alcance.}
\label{tab:alcance-nf}
\begin{tabularx}{\textwidth}{|C{1.6cm}|C{1.7cm}|L{3.5cm}|Y|C{2.1cm}|}
  \hline
  \filaenc \textbf{Clave} & \textbf{Objeto} & \textbf{Característica de calidad} & \textbf{Criterio o umbral acordado} & \textbf{Prioridad} \\ \hline
  \campo{CNF-01} & \campo{OP-02} & \campo{Seguridad — autenticación} & \campo{Umbral verificable} & \campo{Alta} \\ \hline
  & & & & \\ \hline
\end{tabularx}
\notatabla{Una característica no funcional sin criterio o umbral acordado no es verificable: acuérdelo con los interesados antes de diseñar los casos.}
\end{table}

\begin{table}[H]
\centering
\caption{Exclusiones del alcance, su motivo y el riesgo que se asume.}
\label{tab:exclusiones}
\begin{tabularx}{\textwidth}{|C{1.6cm}|L{3.7cm}|Y|C{2.6cm}|}
  \hline
  \filaenc \textbf{Clave} & \textbf{Qué se excluye} & \textbf{Justificación de la exclusión} & \textbf{Riesgo asumido (\ref{sec:riesgos})} \\ \hline
  \campo{EX-01} & \campo{Característica u objeto} & \campo{Por qué no se prueba} & \campo{RP-04} \\ \hline
  & & & \\ \hline
\end{tabularx}
\end{table}

Los Cuadros~\ref{tab:alcance-f} y~\ref{tab:alcance-nf} delimitan lo que sí se probará,
distinguiendo características funcionales de no funcionales; el Cuadro~\ref{tab:exclusiones}
declara lo que queda fuera y, sobre todo, el riesgo que el proyecto acepta al excluirlo. Ese
riesgo debe reaparecer entre los riesgos residuales del apartado~\ref{sec:residuales}.

% ---------------------------------------------------------------------
\subsection{Supuestos y límites}
\label{sec:supuestos}

\begin{instruccion}
Describa los supuestos y limitaciones considerados al elaborar el plan: estándares a cumplir,
política o estrategia de prueba aplicable, tiempo, costo, disponibilidad de personal,
capacitación, herramientas y entornos. Un supuesto es algo que se da por cierto sin haberlo
verificado; si resulta falso, suele convertirse en un riesgo de proyecto.
\end{instruccion}

\begin{table}[H]
\centering
\caption{Supuestos y restricciones del esfuerzo de prueba.}
\label{tab:supuestos}
\begin{tabularx}{\textwidth}{|C{1.6cm}|C{2.2cm}|Y|C{2.6cm}|}
  \hline
  \filaenc \textbf{Clave} & \textbf{Tipo} & \textbf{Enunciado} & \textbf{Si no se cumple} \\ \hline
  \campo{SUP-01} & \campo{Supuesto} & \campo{El entorno estará disponible desde la semana 3} & \campo{RPy-02} \\ \hline
  \campo{LIM-01} & \campo{Restricción} & \campo{No se dispone de datos reales de producción} & \campo{Efecto en el plan} \\ \hline
  & & & \\ \hline
\end{tabularx}
\notatabla{Tipo: \textit{Supuesto} (se da por cierto y puede fallar) · \textit{Restricción} (límite conocido y firme).}
\end{table}

El Cuadro~\ref{tab:supuestos} explicita las condiciones bajo las cuales el plan es válido y
enlaza cada supuesto con el riesgo de proyecto en que se convierte si no se cumple.

\begin{politicaacademica}[Política académica: política general de prueba del curso]
La Experiencia Educativa fija la siguiente política de cobertura mínima \emph{a nivel de
planeación}, en función de la criticidad asignada a cada objeto de prueba
(apartado~\ref{sec:criticidad}):

\medskip
\begin{center}
\begin{tabular}{|L{3.4cm}|C{3.2cm}|}
  \hline
  \filaenc \textbf{Criticidad} & \textbf{Cobertura mínima} \\ \hline
  Catastrófico  & 100\,\% \\ \hline
  Crítico       & 80\,\% \\ \hline
  Marginal      & 60\,\% \\ \hline
  Despreciable  & 40\,\% \\ \hline
\end{tabular}
\end{center}
\medskip

A nivel de ejecución, deberá someterse a prueba la totalidad de los elementos con criticidad
\emph{Catastrófico}. Las demás técnicas abordadas en el curso se aplicarán según corresponda a
ese mismo nivel de integridad, justificando su elección según las características de las
unidades a probar.

\textbf{Estos porcentajes son una regla del curso, no un umbral de \norma.} Para que sean
interpretables debe declararse sobre \emph{qué elemento} se calcula la cobertura: no es lo
mismo cubrir el 80\,\% de los requisitos que el 80\,\% de las particiones de equivalencia o de
las decisiones del código. Defina ese elemento en el apartado~\ref{sec:cobertura}.
\end{politicaacademica}

\begin{instruccion}
Si reutiliza esta plantilla fuera del curso, sustituya la caja anterior por la política de
prueba de su organización, o por umbrales acordados con los interesados y justificados por el
riesgo. Deje constancia de esa adaptación en el apartado~\ref{sec:desviaciones}.
\end{instruccion}

% ---------------------------------------------------------------------
\subsection{Interesados en la prueba}

\begin{instruccion}
Relacione a los interesados (\emph{stakeholders}) en la prueba y su relevancia para el proceso,
indicando los objetivos y expectativas que tienen sobre la prueba en función de su rol.
Identifique también qué información necesita recibir cada uno y con qué periodicidad: ese dato
alimenta directamente el plan de comunicación del apartado~\ref{sec:informes}.
\end{instruccion}

\begin{table}[H]
\centering
\caption{Interesados en la prueba, expectativas y necesidades de información.}
\label{tab:interesados}
\begin{tabularx}{\textwidth}{|L{3.2cm}|L{2.9cm}|Y|C{2.9cm}|}
  \hline
  \filaenc \textbf{Interesado} & \textbf{Rol frente a la prueba} & \textbf{Objetivo o expectativa} & \textbf{Información y frecuencia} \\ \hline
  \campo{Nombre o grupo} & \campo{Decide / informa / provee} & \campo{Qué espera obtener de la prueba} & \campo{Informe de estado semanal} \\ \hline
  & & & \\ \hline
  & & & \\ \hline
\end{tabularx}
\end{table}

El Cuadro~\ref{tab:interesados} identifica a quién sirve el esfuerzo de prueba y qué espera
cada parte; la última columna determina los destinatarios y la periodicidad de los informes de
estado definidos en el apartado~\ref{sec:informes} y formalizados en el \ref{anx:informes}.
